% Document class and font size
\documentclass[a4paper,10pt]{extarticle}

% Compile Command:
% latexmk -pdf resume.tex

% -----------------------------------------------------------
% PACKAGES
% -----------------------------------------------------------
\usepackage{geometry}
\geometry{letterpaper, margin=0.4in}
\usepackage{titlesec}
\usepackage{enumitem}
\usepackage{hyperref}
\usepackage{xcolor}
\usepackage{array}
\usepackage{microtype}
\usepackage{lmodern}

% -----------------------------------------------------------
% FORMATTING
% -----------------------------------------------------------
\setlist{noitemsep}
\hyphenpenalty=10000
\titleformat{\section}{\large\bfseries\uppercase}{\thesection}{1em}{}[\titlerule]
\titlespacing*{\section}{0pt}{8pt}{5pt} % Tight spacing

% Link Formatting
\hypersetup{
    colorlinks=true,
    linkcolor=black,
    filecolor=magenta,
    urlcolor=black,
}

%-------------------------------------------------------------------------------
% BEGIN DOCUMENT
%-------------------------------------------------------------------------------
\begin{document}

% Disable page numbers
\pagestyle{empty}

% Header
\begin{center}
    \textbf{\Huge RUDRA CHAVDA}\\[4pt]
    \small
    \href{mailto:rudrabchavda@gmail.com}{rudrabchavda@gmail.com} \ $|$ \
    \href{https://www.linkedin.com/in/rudra-chavda/}{LinkedIn: rudra-chavda} \ $|$ \
    \href{https://rudra.chavda.online/}{Portfolio: rudra.chavda.online} \ $|$ \
    U.S. Permanent Resident
\end{center}

%------------------------
% Education
\section*{EDUCATION}
    \noindent
    \textbf{University of California, Santa Barbara} \hfill Expected Spring 2028 \\
    \textit{Bachelor of Science, Computer Engineering} \hfill \textbf{GPA: 3.86}

\vspace{5pt}

    \noindent
    \textbf{University of California, Santa Cruz} \hfill 2024 -- 2026 \\
    \textit{Bachelor of Science, Computer Engineering} \hfill \textbf{Honors:} 6x Dean's List

\vspace{2pt}
\noindent
\textbf{Coursework:} Systems programming, computer architecture, operating systems, data structures and algorithms, and software engineering.


%------------------------
% Technical Skills Section
\section*{TECHNICAL SKILLS}
    \noindent
    \begin{tabular}{ @{} >{\bfseries}p{0.12\textwidth} @{\hspace{2ex}} >{\raggedright\arraybackslash}p{0.84\textwidth} @{} }
        Languages & Python, C, C++, Java, TypeScript, JavaScript, SQL, RV Assembly \\
        Frameworks & Next.js, React, Vue.js, Spring Boot, Node.js, Py4Web, Tailwind CSS, Angular \\
        Tools & Git, Docker, Kubernetes, MySQL, Flyway, Microsoft Graph API, JUnit, CARLA, GDAL \\
        Concepts & LLM Agents, Function Calling, OAuth 2.0, Reinforcement Learning, System Design, Data Structures
    \end{tabular}

%------------------------
% Experience
\section*{EXPERIENCE}

    \noindent
    \textbf{Ascend Cargo} \hfill Jun 2026 -- Sep 2026 \\
    \textit{AI/ML Software Engineering Intern}
    \begin{itemize}[leftmargin=0.2in, label={--}]
        \item Architected an enterprise \textbf{LLM agent framework} utilizing \textbf{Java 23, Spring Boot 3.4, and OpenAI function calling}, featuring annotation-driven tool registration and dynamic JSON-Schema generation for targeted intent routing.
        \item Engineered a natural-language shipment management system, diagnosing and resolving a synchronous automation bottleneck that reduced database execution latency from \textbf{12 minutes to under 30 seconds}.
        \item Designed a Role-Based Access Control (RBAC) authorization pipeline enforcing strict per-user AI limits (\textbf{NONE\slash READ\slash WRITE}) and integrated a \textbf{Microsoft Graph API} inbound email channel via \textbf{OAuth 2.0} for authenticated, remote \textbf{shipment creation and updates}.
        \item Hardened agent fault tolerance by implementing explicit refuse-and-instruct error contracts, and authored \textbf{122 JUnit/Mockito unit tests} alongside a custom CLI harness to validate state mutations against the database.
    \end{itemize}
    \vspace{5pt}

    \noindent
    \textbf{AIEA Lab (UCSC)} \hfill Dec 2025 -- Apr 2026 \\
    \textit{Undergraduate Researcher}
    \begin{itemize}[leftmargin=0.2in, label={--}]
        \item Benchmarked and wrote Reinforcement Learning (RL) algorithms including \textbf{DQN, DDQN, A3C, and PPO} to evaluate autonomous vehicle safety parameters within the Robustifying AVs project.
        \item Deployed and scaled the \textbf{CARLA} autonomous driving simulator in a headless environment across the \textbf{Nautilus Kubernetes} hyper-cluster, isolating edge-case failure modes via post-training reward function visualization.
    \end{itemize}
    \vspace{5pt}

    \noindent
    \textbf{CA-Futures (UCSC EEB \& Baskin Engineering)} \hfill May 2025 -- Sep 2025 \\
    \textit{Software Developer}
    \begin{itemize}[leftmargin=0.2in, label={--}]
        \item Engineered a full-stack climate modeling platform utilizing \textbf{Vue.js} and \textbf{Py4Web}, delivering interactive web maps that enable researchers to visualize and compare predictive ecological shifts across multiple long-term climate scenarios.
        \item Architected an asynchronous geospatial pipeline utilizing \textbf{GDAL, Rasterio, and Geopandas} to intercept user-drawn map polygons, compute dynamic habitat connectivity, and implement output caching to eliminate redundant backend execution.
    \end{itemize}
    \vspace{5pt}

    \noindent
    \textbf{Opennote (YC S25)} \hfill Jul 2025 -- Feb 2026 \\
    \textit{Campus Growth Lead}
    \begin{itemize}[leftmargin=0.2in, label={--}]
        \item Spearheaded campus integration strategies, driving a \textbf{30\% increase} in user engagement and securing \textbf{50+ daily active users} by organizing targeted technical workshops on TypeScript user-agent API flows.
        \item Co-directed the UCSC Cruzhacks integration, managing vendor partnerships and allocating \textbf{\$1,500+} in technical resources to drive platform adoption among computer science students.
    \end{itemize}
    \vspace{5pt}

\newpage
%------------------------
% Projects
\section*{PROJECTS}

    \noindent
    \textbf{Better Notes} $|$ \textit{Next.js 15, TypeScript, Convex, Clerk, Tailwind} \hfill Spring 2025
    \begin{itemize}[leftmargin=0.2in, label={--}]
        \item Architected a full-stack, real-time note-taking application, integrating \textbf{Clerk} for secure user authentication and \textbf{Convex} for reactive data synchronization and backend storage.
        \item Implemented a rich-text editing experience utilizing \textbf{BlockNote and ProseMirror}, enabling seamless markdown rendering and dynamic block-based study workflows for end-users.
        \item Scaled application to \textbf{100+ active users} within 3 months of Vercel deployment, optimizing database queries to maintain \textbf{sub-50ms read/write latency}.
    \end{itemize}
    \vspace{5pt}

    \noindent
    \textbf{DivExplorer Bias Detection} $|$ \textit{Python, Pandas, ML, Shapley Values} \hfill Fall 2025
    \begin{itemize}[leftmargin=0.2in, label={--}]
        \item Engineered an automated data analysis pipeline utilizing \textbf{FP-growth and Apriori} frequent pattern mining algorithms in Pandas to extract and evaluate statistically diverging subgroups within large datasets.
        \item Implemented statistical redundancy pruning thresholds and computed \textbf{Shapley values} to quantify specific feature contributions, isolating anomalies and reducing high-overlap duplicate discoveries.
        \item Quantified machine learning classification bias by measuring divergence metrics (e.g., \textbf{false-positive disparities}) across demographic intersections, establishing rigorous baselines for model fairness.
    \end{itemize}
    \vspace{5pt}

    \noindent
    \textbf{Personal Portfolio} $|$ \textit{Next.js 16, TypeScript, Three.js, Motion, Tailwind} \hfill Winter 2024
    \begin{itemize}[leftmargin=0.2in, label={--}]
        \item Developed a high-performance developer portfolio combining Next.js server-side rendering with complex 3D WebGL scenes using \textbf{Three.js} and declarative animations via \textbf{Motion}.
        \item Engineered a custom technical content rendering pipeline, integrating \textbf{KaTeX} for complex mathematical notation and dynamic syntax highlighting for interactive code blocks.
        \item Built a responsive, accessible component architecture utilizing \textbf{Tailwind CSS v4 and Shadcn UI} primitives, optimizing component hydration to maintain a consistent \textbf{60 FPS} experience.
    \end{itemize}

\end{document}